% Fixed131: family-level tame C3 projector, its exact wild-descent ideal,
% a first-order nonextendability witness, and curved two-block transport.
% This subsection is integrated verbatim into the full fixed131 manuscript.

\subsection{The universal tame projector and its wild descent obstruction}
\label{subsec:dyadic-s3-universal-tame-projector}

Version fixed130 proves a canonical decomposition for the
\emph{fixed} tame $S_3$ adjoint lattice.  It does not imply
that the resulting projector is a map of Galois
representations after arbitrary framed deformation:
the wild generators need not preserve its image.
We now separate two questions:
\emph{(a)} whether a canonical $C_3$ projector exists
integrally over the full completed residual ring, and
\emph{(b)} exactly where that projector descends to
an invariant coefficient decomposition.
Both questions admit finite residual-jet formulas.
The answer to (a) is positive; the answer to (b)
is an explicit closed proper locus.

Let
\[
\bar\rho:\Gamma=G_{\mathbf Q_2}
\twoheadrightarrow S_3
\simeq\mathrm{GL}_2(\mathbf F_2)
\]
be the residual representation of
Theorem~\ref{thm:dyadic-s3-integral-sector}.
Let $(A,\mathfrak m,k)$ be a complete
Noetherian local $\mathcal O_E$-algebra,
with a framed lift
$\rho_A:\Gamma\to\mathrm{GL}_2(A)$
whose reduction is $\bar\rho$.
Write
\[
X=\rho_A(\tau),\quad
Y=\rho_A(\sigma),\quad
X_j=\rho_A(x_j)\ (j=0,1),
\qquad
\bar X=\bar U,\quad
\bar Y=\bar S,\quad
\bar X_j=I_2.
\]
The marked tame relation remains
\begin{equation}\label{eq:dyadic-f131-tame-relation}
Y^{-1}XY=X^2.
\end{equation}

\begin{theorem}[Canonical integral order-three extraction in every residual family]
\label{thm:dyadic-f131-canonical-tame-extraction}
Put $Q=X^3-I_2\in M_2(\mathfrak m)$.
The binomial series
\begin{equation}\label{eq:dyadic-f131-tame-root}
u_A:=(X^3)^{1/3}
=\sum_{j\ge0}\binom{1/3}{j}Q^j
\in I_2+M_2(\mathfrak m)
\end{equation}
converges $\mathfrak m$-adically
and is integral over $A$.
The matrix
\begin{equation}\label{eq:dyadic-f131-odd-tame-matrix}
\boxed{\qquad
z_A=Xu_A^{-1}
\qquad}
\end{equation}
is the unique order-three
component of $X$ in its
procyclic closure.  Moreover the
\emph{exact marked tame relation}
forces $u_A=I_2$ and $z_A=X$.
In particular the full universal
rank-two tame matrix already
has exact order three; it satisfies
\begin{equation}\label{eq:dyadic-f131-cubic-etale}
\boxed{\qquad
z_A^3=I_2,\qquad
z_A^2+z_A+I_2=0,\qquad
z_A\bmod\mathfrak m=\bar U.
\qquad}
\end{equation}
The algebra
\begin{equation}\label{eq:dyadic-f131-canonical-tame-etale-algebra}
\mathcal C_A:=A[z_A]\subseteq M_2(A)
\simeq A[T]/(T^2+T+1)
\end{equation}
is a finite étale $A$-algebra
of rank two and is exactly the
centralizer of $z_A$ in $M_2(A)$.

On $M_A=\operatorname{End}_A(A^2)$
define the intrinsic tame Reynolds operator
\begin{equation}\label{eq:dyadic-f131-universal-tame-projector}
\boxed{\qquad
\Pi_A=\frac13
\left(I+\operatorname{Ad}_{z_A}
+\operatorname{Ad}_{z_A^2}\right).
\qquad}
\end{equation}
It is a strict integral idempotent,
with finite projective direct summands
\begin{equation}\label{eq:dyadic-f131-universal-tame-summands}
M_A=M_{0,A}\oplus M_{1,A},
\qquad M_{0,A}=\mathcal C_A=\operatorname{im}\Pi_A,
\qquad M_{1,A}=\ker\Pi_A,
\end{equation}
both of rank two.
This construction commutes with
continuous local coefficient base
change.  On the genuine
marked residual deformation ring
$z_A=X$, so the tame
projector is already a
\emph{finite polynomial} in
$X$ and its adjoint action,
with denominator the unit $3$
only.  The binomial formula
\eqref{eq:dyadic-f131-tame-root}
remains a compatible calculation
of the pro-$2$ part, which is
identically $I_2$ here.
There is no tame power-series
solve, even before passing
to finite precision.

Furthermore the full tame relation
\eqref{eq:dyadic-f131-tame-relation}
forces
\begin{equation}\label{eq:dyadic-f131-sigma-tame-normalization}
\boxed{\qquad
Y^{-1}z_A Y=z_A^2,
\quad
[\operatorname{Ad}_Y,\Pi_A]=0,
\quad
[\operatorname{Ad}_X,\Pi_A]=0.
\qquad}
\end{equation}
Thus the potential failure of
Galois-equivariant splitting comes
\emph{only} from wild transport.
\end{theorem}

\begin{proof}
The matrix $Q=X^3-I_2$ lies in
$M_2(\mathfrak m)$ because the
residual tame matrix has order
three.  The coefficients
$\binom{1/3}{j}\in\mathbf Z_2$
are integral and the $j$th
summand belongs to
$M_2(\mathfrak m^j)$.
The usual binomial identity,
valid for powers of the single
commuting matrix $Q$, proves
$u_A^3=X^3$.
Since $u_A$ is a power series
in $X$, it commutes with $X$,
so $z_A^3=I_2$.
Its reduction is $\bar U$.
The matrix $z_A-I_2$ is invertible:
its reduction $\bar U-I_2$
has determinant one.
Multiplying
$(z_A-I_2)(z_A^2+z_A+I_2)=0$
by its inverse proves the
quadratic polynomial in
\eqref{eq:dyadic-f131-cubic-etale}.

The polynomial $T^2+T+1$
is separable in residue
characteristic two; its
derivative $2T+1$ is a unit.
The map
$A[T]/(T^2+T+1)\to M_2(A)$
sending $T$ to $z_A$
is injective, since
$I_2,z_A$ are linearly
independent after reduction.
It is a finite étale rank-two
algebra.
More concretely, for the
first framed basis vector $e_1$,
the pair $(e_1,z_A e_1)$
is an $A$-basis because its
reduction is a basis for
$\bar U$.  In this cyclic basis,
every matrix commuting with $z_A$
is a unique $A$-linear
polynomial of degree at most
one in $z_A$.
This proves the centralizer
claim and its rank.

The conjugation
$\operatorname{Ad}_{z_A}$
has order three, so its
Reynolds average is an
idempotent over $A$.
Its image is precisely the
centralizer $\mathcal C_A$.
Its reduction is the rank-two
idempotent of
Theorem~\ref{thm:dyadic-s3-adjoint-unimodular-splitting};
the complementary summand is
therefore finite free of rank
two by Nakayama and local
projectivity.
All constructions use integral
convergent power series and
unit matrices and hence commute
with completed local base
change.  Terms of order
$j\ge t$ vanish modulo
$\mathfrak m^t$.

Finally the exact tame relation
also proves the stronger
\emph{rigidity} claim.
Let $d=\det X\in A^\times$.
Conjugating and taking
determinants in
\eqref{eq:dyadic-f131-tame-relation}
gives $d=d^2$, hence $d=1$.
Put $t=\operatorname{tr}X$.
The two sides of the tame
relation have equal trace,
so Cayley--Hamilton yields
\[
t=\operatorname{tr}(X^2)=t^2-2d=t^2-2,
\qquad (t-2)(t+1)=0.
\]
Since $\bar t=\operatorname{tr}(\bar U)=1$
in residue characteristic two,
$t-2$ is a unit.
Therefore $t=-1$.
Cayley--Hamilton now gives
$X^2+X+I_2=0$ and $X^3=I_2$.
Consequently the unique
pro-$2$ cubic root of $X^3$
in \eqref{eq:dyadic-f131-tame-root}
is $u_A=I_2$, and
$z_A=X$ identically.
In particular the claimed
tame projector is \emph{finite
polynomial} on the full
marked deformation ring.
The tame relation directly gives
$Y^{-1}z_A Y=z_A^2$.
Conjugation by $Y$ therefore
normalizes the cyclic
order-three adjoint action,
and commutes with its
Reynolds projector.
The matrix $X$ commutes with
its own odd component $z_A$,
so its adjoint action does as
well.  This proves all
assertions.
\end{proof}

\begin{theorem}[A fully native universal tame frame in four transport matrices]
\label{thm:dyadic-f131-native-tame-four-frame}
For any complete framed
$S_3$ residual deformation
over $A$, let
$X=\rho_A(\tau)$ and
$Y=\rho_A(\sigma)$.
The tame rigidity theorem
gives
\[
X^2+X+I_2=0,
\qquad Y^{-1}XY=X^2.
\]
Then the four matrices
\begin{equation}\label{eq:dyadic-f131-universal-native-frame}
\boxed{\qquad
\mathcal B_A=(I_2,X,Y,XY)
\qquad}
\end{equation}
form an \emph{integral
$A$-basis of the entire}
rank-four algebra
$M_2(A)$.
Indeed the matrix with
their row-major vectorizations
as columns reduces to
\begin{equation}\label{eq:dyadic-f131-native-frame-residue}
\mathsf B_{\rm res}=
\begin{pmatrix}
1&-1&0&-1\\
0&-1&1&-1\\
0&1&1&0\\
1&0&0&1
\end{pmatrix},
\qquad
\boxed{\det\mathsf B_{\rm res}=-3
\in\mathbf Z_2^\times}.
\end{equation}
In this \emph{transport-native}
basis the tame projector
has the exact global
matrix, for \emph{every}
framed deformation:
\begin{equation}\label{eq:dyadic-f131-universal-diagonal-projector}
\boxed{\qquad
[\Pi_A]_{\mathcal B_A}
=\operatorname{diag}(1,1,0,0),
\qquad
M_{0,A}=A I_2\oplus AX,
\quad
M_{1,A}=AY\oplus AXY.
\qquad}
\end{equation}
Moreover both tame
generator actions are
strictly block diagonal,
with their complete
integral matrices
\begin{equation}\label{eq:dyadic-f131-universal-block-actions}
\boxed{
\begin{aligned}
[\operatorname{Ad}_X]_{\mathcal B_A}
&=
\operatorname{diag}\left(
 I_2,\
 \begin{pmatrix}-1&1\\-1&0\end{pmatrix}
 \right),\\[2pt]
[\operatorname{Ad}_Y]_{\mathcal B_A}
&=
\operatorname{diag}\left(
 \begin{pmatrix}1&-1\\0&-1\end{pmatrix},
 \begin{pmatrix}1&-1\\0&-1\end{pmatrix}
 \right).
\end{aligned}}
\end{equation}
In particular
\begin{equation}\label{eq:dyadic-f131-frobenius-square-scalar}
\boxed{\qquad Y^2\in A^\times I_2. \qquad}
\end{equation}
No matrix logarithm,
period-module chart,
universal Fox syzygy,
or binomial power
extraction is needed
to construct this
four-dimensional
native tame frame.

For $j=0,1$ the wild
generator matrix has
a unique expansion
\begin{equation}\label{eq:dyadic-f131-wild-native-components}
\boxed{
X_j=a_j I_2+b_j X
+\eta_jY+\xi_jXY,
\qquad
a_j,b_j,\eta_j,\xi_j\in A.
}
\end{equation}
The \emph{exact} split
ideal is therefore
given by the four
native scalar
equations
\begin{equation}\label{eq:dyadic-f131-four-scalar-split-ideal}
\boxed{\qquad
\mathfrak I_{\rm split}
=(\eta_0,\xi_0,\eta_1,\xi_1)
\subset R_{\bar\rho}^{\square}.
\qquad}
\end{equation}
They are finite
matrix/minor
expressions on
the complete
residual chart.
Only two independent
linear equations
remain at the
residual tangent
point, as quantified
in
Corollary~\ref{cor:dyadic-f131-split-tangent-codimension}.
No global regular-sequence
claim is implicit.
\end{theorem}

\begin{proof}
At the residual
$S_3$ matrices
$U,S$, the
row-major matrix
with columns
$\operatorname{vec}(I_2)$,
$\operatorname{vec}(U)$,
$\operatorname{vec}(S)$,
$\operatorname{vec}(US)$
is
\eqref{eq:dyadic-f131-native-frame-residue}
and has determinant
$-3$, an odd unit.
Thus its lift
$\mathsf B_A$
is invertible
over every complete
local $A$ with
the chosen residual
marking.
The first two
basis matrices
commute with
$X$ and belong
to $\operatorname{im}\Pi_A$.
From $Y^{-1}XY=X^2$
one obtains
\[
XY=YX^2,\qquad
YX=X^2Y,\qquad
XYX^{-1}=X^2Y.
\]
Hence conjugation
by $X$ acts on the
second pair by
\[
\operatorname{Ad}_X(Y)
=X^2Y=-Y-XY,\qquad
\operatorname{Ad}_X(XY)=Y.
\]
These equations
give the first
matrix of
\eqref{eq:dyadic-f131-universal-block-actions}.
Its second block
has order three
and satisfies
$I+C+C^2=0$.
Therefore the
Reynolds average
vanishes on
$AY\oplus AXY$
and is identity
on $AI_2\oplus AX$,
proving the
diagonal projector.

The normalizer
action of $Y$ on
the rank-two
algebra $A[X]$
has order two;
the tame relation
and $X^3=1$
give
$YXY^{-1}=X^2$.
Thus
\[
\operatorname{Ad}_Y(X)
=-I_2-X,\qquad
\operatorname{Ad}_Y(Y)=Y,\qquad
\operatorname{Ad}_Y(XY)
=X^2Y=-Y-XY,
\]
proving the
second matrix in
\eqref{eq:dyadic-f131-universal-block-actions}.
Its square is the
identity on the
\emph{entire}
basis of $M_2(A)$.
Hence
$\operatorname{Ad}_{Y^2}$
is trivial on
$M_2(A)$, and
its centralizer
consists of scalars.
Since $Y$ is
invertible,
$Y^2$ is a
scalar unit.

Finally the
basis gives
the unique
expansion of
each wild matrix.
By the diagonal
projector just proved,
\[
(1-\Pi_A)X_j
=\eta_jY+\xi_jXY.
\]
Therefore the
matrix centralizer
equations of
Theorem~\ref{thm:dyadic-f131-universal-split-locus}
are equivalent
to $\eta_j=\xi_j=0$
for $j=0,1$,
proving the
explicit four-equation
ideal formula.
Every coordinate
is recovered by
inverting the
unit determinant
$\det\mathsf B_A$,
so it is an
integral finite
matrix expression
on the whole
completed residual
chart.
\end{proof}

\begin{theorem}[Exact wild equations for the universal split locus]
\label{thm:dyadic-f131-universal-split-locus}
Let $R_{\bar\rho}^{\square}$ be
the framed complete local
Galois deformation ring
and form $z^\square,\Pi^\square$
using Theorem~\ref{thm:dyadic-f131-canonical-tame-extraction}.
Let $\mathfrak I_{\rm split}$
be the closed ideal of
$R_{\bar\rho}^{\square}$
generated by the matrix entries of
\begin{equation}\label{eq:dyadic-f131-wild-split-equations}
\boxed{\qquad
[X_0,z^\square],\qquad
[X_1,z^\square].
\qquad}
\end{equation}
Equivalently, by
\eqref{eq:dyadic-f131-universal-tame-summands},
the ideal is generated by
the two rank-two block
components
\begin{equation}\label{eq:dyadic-f131-wild-split-projector-equations}
\boxed{\qquad
(1-\Pi^\square)X_0,\qquad
(1-\Pi^\square)X_1,
\qquad}
\end{equation}
where $\Pi^\square$ acts
on the \emph{matrix} algebra.
On any framed deformation
over a complete local $A$,
the following are equivalent:
\begin{enumerate}[label=\textnormal{(\roman*)}]
\item $\operatorname{Ad}_{\rho_A(g)}$
commutes with $\Pi_A$
for every $g\in\Gamma$;
\item both blocks $M_{0,A}$ and
$M_{1,A}$ of
\eqref{eq:dyadic-f131-universal-tame-summands}
are $\Gamma$-subrepresentations;
\item $X_0,X_1$ commute with $z_A$;
\item every marked generator
normalizes the finite étale
algebra $\mathcal C_A$,
hence
\[
\rho_A(\Gamma)\subseteq
N_{\mathrm{GL}_2(A)}(\mathcal C_A).
\]
\end{enumerate}
Therefore
\begin{equation}\label{eq:dyadic-f131-split-ring}
R_{\rm split}:=
R_{\bar\rho}^{\square}/\mathfrak I_{\rm split}
\end{equation}
represents the \emph{exact}
tame-isotypic split
subfunctor of the framed
residual deformation functor.

Over $R_{\rm split}$ there
is an integral Galois-equivariant
coefficient decomposition,
stable under all derived
base changes within the
split subfunctor:
\begin{equation}\label{eq:dyadic-f131-split-cochains}
R\Gamma(\Gamma,M_{R_{\rm split}})
\simeq
R\Gamma(\Gamma,M_{0,R_{\rm split}})
\oplus
R\Gamma(\Gamma,M_{1,R_{\rm split}}).
\end{equation}
The same splitting applies
to the all-Sylow finite
perfect coefficient carrier.
This does \emph{not} assert
the fixed130 constant Smith
normal form on the entire
split family.

The étale algebra
$\mathcal C_A$ makes $A^2$
a free rank-one
$\mathcal C_A$-module.
On the split locus the
index-two Galois subgroup
$G_{\mathbf Q_2(\zeta_3)}$
acts $\mathcal C_A$-linearly
and the other coset
acts through the
nontrivial involution
$z_A\mapsto z_A^2$.
Thus the split locus is
precisely the associated
rank-one étale-toric/
semilinear normalizer
realization, not an
arbitrary extension of
the fixed $S_3$ matrix frame.
\end{theorem}

\begin{proof}
The equivalence of (i)
and (ii) is idempotent
linear algebra.
By
\eqref{eq:dyadic-f131-sigma-tame-normalization},
$\sigma$ and $\tau$
already commute with
$\Pi_A$ on the adjoint
module.
Suppose that an individual
wild matrix $X_j$, whose
residue is $I_2$, satisfies
$[\operatorname{Ad}_{X_j},\Pi_A]=0$.
Then $\operatorname{Ad}_{X_j}$
preserves
$\mathcal C_A=\operatorname{im}\Pi_A$,
so $X_j z_A X_j^{-1}$
lies in $\mathcal C_A$,
is a cubic root of one,
and reduces to $\bar U$.
As $3z_A^2$ is a unit
in the commutative
$\mathfrak m$-complete
algebra $\mathcal C_A$,
the root lifting the given
residue is unique.
Consequently
$X_jz_AX_j^{-1}=z_A$.
Conversely, if $X_j$
commutes with $z_A$,
its adjoint action
commutes with the
Reynolds projector.
This proves (i)--(iii).

Every marked generator
normalizes $\mathcal C_A$
if (iii) holds, because
$\sigma$ acts on it by
$z_A\mapsto z_A^2$,
while $\tau,x_0,x_1$
centralize it.
Conversely, if a wild
$X_j$ normalizes
$\mathcal C_A$, the
same unique-root argument
shows it centralizes $z_A$.
This proves (iv) and the
functorial vanishing
criterion.
Since all equations are
entries of matrices over
the complete universal
deformation ring,
they define the stated
closed represented subfunctor.
The equivalent projected
form follows from
$\ker\operatorname{ad}(z_A)
=\mathcal C_A=\operatorname{im}\Pi_A$.

Over that quotient the two
idempotent summands are
preserved by $\Gamma$.
Ordinary continuous cochains
and the perfect all-Sylow
carrier are additive in
the coefficient module,
giving the displayed
derived splitting.
Its base-change
compatibility follows
because the idempotent and
all actions are defined
already over $R_{\rm split}$.

The cyclic basis
$(e_1,z_Ae_1)$ makes
$A^2$ a free rank-one
module over $\mathcal C_A$.
The normalizer acts either
linearly over $\mathcal C_A$
or semilinearly under its
unique nontrivial étale
involution.  The tame
generator $\sigma$ acts
semilinearly, while the
even-parity index-two
subgroup is generated
inside the normalizer
by $\tau$, the wild
generators and
$\sigma^2$ and therefore
acts linearly.
This proves the final
description without
claiming a global
minimal Fox basis.
\end{proof}

\begin{theorem}[A semilinear rank-one description of the entire split subfunctor]
\label{thm:dyadic-f131-toric-semilinear-description}
Keep the complete split
deformation ring
$R_{\rm split}$ of
Theorem~\ref{thm:dyadic-f131-universal-split-locus}
and let $A$ be any
complete local point of it.
Put
\[
\Gamma'=
G_{\mathbf Q_2(\zeta_3)}
\triangleleft\Gamma,\qquad
\mathcal C=A[z_A],\qquad
\iota(z_A)=z_A^2.
\]
Fix the framed cyclic
vector $e_1\in A^2$.
The map
\[
\mathcal C\longrightarrow A^2,
\qquad a\longmapsto a e_1
\]
is an isomorphism
of rank-one
$\mathcal C$-modules.
Under this identification
the restriction of $\rho_A$
to $\Gamma'$ is exactly
multiplication by a
continuous character
\begin{equation}\label{eq:dyadic-f131-toric-character}
\chi:\Gamma'\longrightarrow\mathcal C^\times,
\qquad
\rho_A(h)=m_{\chi(h)}
\quad(h\in\Gamma').
\end{equation}
The other Galois coset
has the explicit form
\begin{equation}\label{eq:dyadic-f131-semilinear-sigma}
\rho_A(\sigma)=m_b\circ\iota,
\qquad b\in\mathcal C^\times.
\end{equation}
Conversely, rank-one
$\mathcal C$-valued
character data $(\chi,b)$
give a split framed
representation precisely
when their residual
marking agrees with
$\bar\rho$ and
\begin{equation}\label{eq:dyadic-f131-semilinear-character-conditions}
\boxed{
\begin{aligned}
\chi(\sigma h\sigma^{-1})
 &=\iota(\chi(h))
 &&(h\in\Gamma'),\\
b\,\iota(b)&=\chi(\sigma^2).
\end{aligned}}
\end{equation}
Thus the split locus
is a genuine
\emph{quadratic-étale
rank-one/semilinear
descent locus},
not merely a
set-theoretic
centralizer condition.

Moreover the two
adjoint summands
have the exact
integral realizations
\begin{equation}\label{eq:dyadic-f131-skew-endo-split}
\boxed{\quad
\operatorname{End}_A(A^2)
=
\underbrace{\{m_a:a\in\mathcal C\}}_{M_{0,A}}
\ \oplus\
\underbrace{\{m_c\circ\iota:c\in\mathcal C\}}_{M_{1,A}},
\quad}
\end{equation}
and the conjugation
action on them is:
\begin{equation}\label{eq:dyadic-f131-skew-endo-actions}
\boxed{
\begin{array}{c|cc}
& h\in\Gamma'&\sigma\\ \hline
M_0:\ a
 &a&\iota(a)\\
M_1:\ c
 &(\chi(h)/\iota\chi(h))c&
 (b/\iota(b))\,\iota(c).
\end{array}}
\end{equation}
In particular
$M_{0,A}$ has
the \emph{same}
integral permutation
action as in fixed130,
for \emph{every}
split deformation:
\begin{equation}\label{eq:dyadic-f131-family-shapiro-m0}
\boxed{\quad
M_{0,A}
\simeq\operatorname{Ind}_{\Gamma'}^\Gamma A,
\qquad
R\Gamma(\Gamma,M_{0,A})
\simeq R\Gamma(\Gamma',A).
\quad}
\end{equation}
This is a derived
Shapiro isomorphism
over the complete
coefficient algebra $A$;
it does not require
the cohomology of $A$
to have constant
rank or to be
$A$-flat.
The second summand
has a generally
\emph{nontrivial}
rank-one
$\mathcal C$-valued
wild character
$\chi/\iota\chi$.
Consequently it
need not be a
constant natural-$T$
summand away from
the original
fixed tame point.
\end{theorem}

\begin{proof}
The cyclic-basis
argument of
Theorem~\ref{thm:dyadic-f131-canonical-tame-extraction}
makes $A^2$ a free
rank-one $\mathcal C$-module.
On the split locus
the full image
normalizes $\mathcal C$.
The kernel of its
residual involution
character is
$\Gamma'$, so
each $h\in\Gamma'$
acts $\mathcal C$-linearly
and is multiplication
by a unique unit
$\chi(h)$.  The map
$\chi$ is continuous
and multiplicative.
The marked generator
$\sigma$ acts by the
nontrivial étale
involution, hence
has the form
$m_b\circ\iota$
for an invertible
coefficient $b$.
The displayed
relations are exactly
the conjugation
and square identities
for this index-two
group extension.
They suffice to
reconstruct a continuous
representation of
$\Gamma$ from
$\chi$ and $b$,
because every element
has a unique
expression in
$\Gamma'$ or
$\sigma\Gamma'$.
The framed residual
condition picks out
the prescribed
component.

For a quadratic
étale Galois
algebra $\mathcal C/A$,
the two families
of $A$-linear
operators $m_a$
and $m_c\circ\iota$
span and are
independent in
$\operatorname{End}_A(\mathcal C)$.
This can be checked
after faithfully
flat étale base
change splitting
$\mathcal C\simeq A\times A$,
where they are
the diagonal and
off-diagonal
matrix blocks,
and hence descends.
The conjugation
by $z_A$ acts
trivially on
$m_a$ and by
multiplication
$z_A/\iota(z_A)=z_A^2$
on $m_c\iota$.
The tame average
therefore fixes
the former and
kills the latter,
giving
\eqref{eq:dyadic-f131-skew-endo-split}.

Conjugating these
two kinds of
operators by
$m_{\chi(h)}$
and $m_b\iota$
gives the four
exact formulas
\eqref{eq:dyadic-f131-skew-endo-actions}.
On the first
summand $\Gamma'$
acts trivially,
and $\sigma$
interchanges the
two basis vectors
$z_A,z_A^2$.
Thus as a
$\Gamma$-module
$M_{0,A}$ is the
permutation-induced
rank-two module
from the trivial
$A$-module on
$\Gamma'$, proving
\eqref{eq:dyadic-f131-family-shapiro-m0}
by integral
Shapiro's lemma.
No averaging by
two or constant
cohomology-rank
assumption occurs.
\end{proof}

\begin{example}[A nontrivial integral wild point on the split locus]
\label{ex:dyadic-f131-nontrivial-split-wild-point}
Over $\mathbf Z_2$ keep
the fixed matrices
$S,U$ of
\eqref{eq:dyadic-s3-integral-matrices}
and assign
\[
\rho(\sigma)=S,\quad
\rho(\tau)=U,\quad
\rho(x_0)=-I_2,\quad
\rho(x_1)=I_2.
\]
The wild generator
$x_0$ has nontrivial
order-two image,
but it is central
and hence commutes
with the tame
order-three matrix.
Indeed
$(x_0\tau)^{\omega_2}=-I_2$,
$d_0=1$, and the
literal marked
wild relation
reduces to
$x_0^2=1$.
The tame relation
holds because
$S^{-1}US=U^2$.
The wild normal
closure is the
central $2$-group
$\{\pm I_2\}$,
so the marked
pro-$2$ condition
is satisfied.
This therefore
defines an
actual integral
Galois representation
on the split
subfunctor, different
from the original
finite $S_3$
representation.
\end{example}

\begin{theorem}[A genuine wild first-order obstruction and failure of universal splitting]
\label{thm:dyadic-f131-first-order-nosplit}
Let
$k=\mathbf F_2$ and
$A_\epsilon=k[\epsilon]/(\epsilon^2)$.
Reduce the fixed matrices
$S,U,V$ of
Theorem~\ref{thm:dyadic-s3-adjoint-unimodular-splitting}
modulo two:
\begin{equation}\label{eq:dyadic-f131-wild-counterexample-matrices}
\bar S=
\begin{pmatrix}0&1\\1&0\end{pmatrix},
\quad
\bar U=
\begin{pmatrix}1&1\\1&0\end{pmatrix},
\quad
\bar V=
\begin{pmatrix}1&1\\0&1\end{pmatrix}.
\end{equation}
The marked generator assignment
\begin{equation}\label{eq:dyadic-f131-wild-counterexample-rep}
\boxed{\quad
\rho_\epsilon(\sigma)=\bar S,\quad
\rho_\epsilon(\tau)=\bar U,\quad
\rho_\epsilon(x_0)=I_2+\epsilon\bar V,\quad
\rho_\epsilon(x_1)=I_2
\quad}
\end{equation}
satisfies the literal
Roe--Turturean tame and
wild relations, including
the marked pro-$2$ wild
normal-closure condition.
It therefore defines a
genuine continuous
$G_{\mathbf Q_2}$ residual
deformation.

Its canonical odd tame
component is
$z_\epsilon=\bar U$.
However
\begin{equation}\label{eq:dyadic-f131-nonzero-wild-commutator}
\boxed{\quad
[\bar V,\bar U]=\bar V\ne0,\qquad
[\rho_\epsilon(x_0),z_\epsilon]
=\epsilon\bar V\ne0.
\quad}
\end{equation}
Thus $\rho_\epsilon$ lies
outside the split locus of
Theorem~\ref{thm:dyadic-f131-universal-split-locus}.

Moreover this failure
cannot be repaired by
replacing the canonical
projector with \emph{any}
Galois-equivariant
idempotent
$E_\epsilon$ that reduces
to the fixed130 tame
projector $\Pi_0$.
Hence the fixed130
$\operatorname{ad}T=M_0\oplus M_1$
splitting does \emph{not}
extend to the unrestricted
universal framed
$S_3$ deformation space,
even after deforming the
projector itself.
\end{theorem}

\begin{proof}
Over the dual-number ring
one has
$\bar U^3=I_2$ and
$\bar S^{-1}\bar U\bar S
=\bar U^2$.
Also
\[
\bar V+\operatorname{Ad}_{\bar U}\bar V
+\operatorname{Ad}_{\bar U^2}\bar V=0.
\]
Therefore
\[
\bigl((I_2+\epsilon\bar V)\bar U\bigr)^3
=I_2
\]
and the $\omega_2$
component of
$\rho_\epsilon(x_0\tau)$
is the identity.
Similarly the
$\omega_2$ component
of $\bar U$ is the
identity, whereas
$\bar S^{\omega_2}=\bar S$.
Consequently the literal
marked auxiliary words
satisfy
\[
u_0=u_1=1,\quad
d_0=x_0^{-1}=x_0,\quad
g_0=\sigma^2=1,\quad
c_0=h_c=1.
\]
Because $x_0^2=1$ and
$d_0=x_0$, the remaining
word $h_0$ is also one.
The full wild relator
$h_0u_1^{-1}x_1^\sigma c_0$
is therefore one, and the
tame relator was checked
above.
The wild images lie in the
normal subgroup
$1+\epsilon M_2(k)$,
which is an elementary
abelian $2$-group.
The marked pro-$2$
condition is automatic.
The universal property
of the marked
Roe--Turturean presentation
then gives an actual
continuous Galois
representation.

Since $\rho_\epsilon(\tau)$
has exact order three,
its pro-$2$ component
is one and
$z_\epsilon=\bar U$.
Direct characteristic-two
matrix multiplication gives
$[\bar V,\bar U]=\bar V$,
proving
\eqref{eq:dyadic-f131-nonzero-wild-commutator}.

Suppose that
$E_\epsilon=\Pi_0+\epsilon D$
were a Galois-equivariant
idempotent reducing to
$\Pi_0$.
As an action on the
adjoint module,
\[
\operatorname{Ad}_{\rho_\epsilon(x_0)}
=I+\epsilon\operatorname{ad}_{\bar V}.
\]
Commutation with $E_\epsilon$
would imply, already
at order $\epsilon$,
\[
[\operatorname{ad}_{\bar V},\Pi_0]=0.
\]
But $\Pi_0(\bar U)=\bar U$,
$\operatorname{ad}_{\bar V}(\bar U)=\bar V$,
and $\Pi_0(\bar V)=0$,
so this commutator maps
$\bar U$ to the nonzero
matrix $\bar V$.
This contradiction is
independent of $D$ and
proves the strong
nonextendability.
\end{proof}

\begin{corollary}[Exact first-order codimension of the split subfunctor]
\label{cor:dyadic-f131-split-tangent-codimension}
At the residual $S_3$
point over $\mathbf F_2$,
the first-order split
condition on a framed
Galois cocycle $c$ is
\begin{equation}\label{eq:dyadic-f131-tangent-wild-character}
\boxed{\quad
\kappa_{\rm wild}([c])
:=(I-\bar\Pi_0)c(x_0)=0
\quad\text{in }\bar M_1.
\quad}
\end{equation}
The map
\[
\kappa_{\rm wild}:
H^1(\Gamma,\operatorname{ad}\bar T)
\longrightarrow\bar M_1
\cong\mathbf F_2^2
\]
is a \emph{surjection}
with kernel
$H^1(\Gamma,\bar M_0)$.
Thus there is a short
exact sequence
\begin{equation}\label{eq:dyadic-f131-tangent-exact-sequence}
\boxed{
0\to H^1(\Gamma,\bar M_0)
\longrightarrow
H^1(\Gamma,\operatorname{ad}\bar T)
\xrightarrow{\;\kappa_{\rm wild}\;}
\bar M_1\to0.
}
\end{equation}
The unframed residual
tangent dimensions are
\[
\dim_{\mathbf F_2}H^1(\Gamma,\operatorname{ad}\bar T)=6,
\qquad
\dim_{\mathbf F_2}\ker\kappa_{\rm wild}=4.
\]
The corresponding framed
tangent dimensions are
$9$ and $7$, respectively.
Hence the tame-isotypic
split subfunctor is a
\emph{proper} closed
subfunctor, with tangent
codimension exactly two.
No assertion that its
global defining ideal is
a regular sequence is made.
\end{corollary}

\begin{proof}
At the residual point
the wild generators
act trivially, so
$c(x_j)$ is unchanged
under infinitesimal
conjugation.
Differentiating the
commutator equations
\eqref{eq:dyadic-f131-wild-split-equations}
gives
$[c(x_j),\bar U]=0$,
equivalently
$(I-\bar\Pi_0)c(x_j)=0$.
The fixed130 blockwise
linearized relations show
that the $\bar M_1$
component of $c(x_1)$
is already zero for
\emph{every} global
cocycle, whereas the
$\bar M_1$ component of
$c(x_0)$ is unrestricted.
Indeed the $\bar M_1$
wild differential reduces
to $\bar S^{-1}c(x_1)$
in characteristic two.
The norm-zero natural
block of fixed129
identifies its $H^1$
with $\bar M_1$ by
evaluation at $x_0$.
Consequently
$\kappa_{\rm wild}$
is the canonical
projection onto the
$\bar M_1$ summand of
the fixed130
decomposition, giving
the short exact
sequence.

The fixed130 $M_0$
Fox complex has
$H^1$ of dimension
four after mod-two
reduction: its three
integral free directions
and its one top-degree
$2$-torsion Bockstein
direction.
The $\bar M_1$
complex contributes two
additional dimensions.
Thus the full unframed
$H^1$ dimension is six
and the kernel has
dimension four.
The residual adjoint
invariants are scalars
of dimension one, so
the degree-zero
coboundary image has
dimension $4-1=3$.
Adding these gauge
directions gives the
framed tangent
dimensions $9$ and $7$.
The explicit deformation
of
Theorem~\ref{thm:dyadic-f131-first-order-nosplit}
realizes a nonzero
obstruction class, so
the ideal is genuinely
proper.
\end{proof}

\begin{theorem}[Exact curved two-block Cartan multiplication]
\label{thm:dyadic-f131-curved-tame-transport}
For an arbitrary framed
deformation over $A$,
retain its canonical tame
projector $\Pi=\Pi_A$
and put
$P_0=\Pi$, $P_1=I-\Pi$.
For $g\in\Gamma$ define
\[
\mathcal R(g)=\operatorname{Ad}_{\rho_A(g)}
\quad\text{on }M_A,
\qquad
\mathcal R_{ab}(g)=P_a\mathcal R(g)P_b,
\quad a,b\in\{0,1\}.
\]
Write
\begin{equation}\label{eq:dyadic-f131-curved-two-blocks}
\mathcal R(g)=
\begin{pmatrix}
A_g&B_g\\ C_g&D_g
\end{pmatrix}
\quad\text{on }M_{0,A}\oplus M_{1,A}.
\end{equation}
The true profinite
transport remains
multiplicatively flat,
but its two truncated
diagonal carriers obey
the \emph{exact} formulas
\begin{equation}\label{eq:dyadic-f131-curved-block-products}
\boxed{
\begin{aligned}
A_{gh}&=A_gA_h+B_gC_h,\\
B_{gh}&=A_gB_h+B_gD_h,\\
C_{gh}&=C_gA_h+D_gC_h,\\
D_{gh}&=C_gB_h+D_gD_h.
\end{aligned}}
\end{equation}
In particular
\begin{equation}\label{eq:dyadic-f131-two-step-return-curvature}
\boxed{\quad
\mathscr F_{00}(g,h)
:=A_{gh}-A_gA_h=B_gC_h,
\quad
\mathscr F_{11}(g,h)
:=D_{gh}-D_gD_h=C_gB_h.
\quad}
\end{equation}
Both diagonal defects
vanish on the split
locus, but outside it
the off-diagonal
wild transport is a
genuine first-order
source of possible
two-step return defects.
The off-diagonal blocks
vanish identically for
$\sigma$ and $\tau$,
and all their nonzero
sources are carried by
wild words.

For any finite residual
precision $\mathfrak m^t$,
each block matrix,
its product formula,
and the projectors
are finite computations
from the marked
transport matrices
and unit-matrix inverses.
The tame projector itself
is already polynomial
on this marked residual
sector, without any
binomial truncation;
only other marked
profinite word evaluations
may require the fixed126
finite-jet procedure.
These formulas define
a \emph{source-faithful
curved two-block
coefficient chart} on
the unrestricted
$S_3$ residual family.
They are a
finite arithmetic
transport identity,
not an assertion of
additional analytic
differential equations.
\end{theorem}

\begin{proof}
The group law is
$\mathcal R(gh)
=\mathcal R(g)\mathcal R(h)$.
Insert $P_0+P_1=I$
between the two factors
and project to each of
the four pairs of block
degrees.  This gives
\eqref{eq:dyadic-f131-curved-block-products}
and
\eqref{eq:dyadic-f131-two-step-return-curvature}
exactly, with no
approximation.
By
\eqref{eq:dyadic-f131-sigma-tame-normalization},
the $\sigma,\tau$
transport preserves
both tame summands.
Their off-diagonal
blocks therefore vanish.
Vanishing of all wild
off-diagonal blocks is
equivalent to the split
criterion of
Theorem~\ref{thm:dyadic-f131-universal-split-locus}.
The stated finite-jet
effectivity follows
from the binomial
tail estimate and
ordinary finite
matrix multiplication.
\end{proof}

\begin{remark}[Two distinct meanings of closure]
\label{rem:dyadic-f131-split-versus-global-atlas}
The universal tame
\emph{matrix projector}
exists over the full
residual deformation
ring.  A
\emph{Galois-equivariant
decomposition} by that
projector exists
precisely on the
wild-centralizer
subfunctor, and a
genuine first-order
counterexample shows
that the latter is
proper.
By contrast the
all-Sylow perfect
cochain carrier and
the Cayley--Hamilton
full-transport
realization exist
over the whole
framed residual
sector regardless
of this failure.
Therefore failure of
a strict tame
eigensplitting does
not reopen the
derived faithful
finite-atlas theorem.
Nor does the curved
block identity
provide the still
missing canonical
Schreier-to-minimal-
Demu\v{s}kin Fox
syzygy comparison
for arbitrary
non-pro-$2$
deformation families.
\end{remark}
